Após executar o experimento aplicando um degrau de amplitude
$50^\circ\mathrm{C}$ como entrada identificamos o comportamento dinâmico da
planta. Para a identificação foi utilizada a abordagem de Ziegler-Nichols por
ensaio ao degrau, adotando um modelo \gls{fopdt}, com isso obtivemos duas
aproximações como plantas, uma para cada sensor, de acordo com a
\Eqref{planta-sensor-1-zn,planta-sensor-2-zn}:

\begin{align}
\eqlabel{planta-sensor-1-zn}
G_{1}(s) &= \frac{0.92}{188.03s + 1} \cdot e^{-6.37s}\\
\eqlabel{planta-sensor-2-zn}
G_{2}(s) &= \frac{0.91}{183.94s + 1} \cdot e^{-11.66s}
\end{align}